\textbf{Phần 1. Trắc nghiệm lựa chọn}
\setcounterpageref{ex}{0}

\begin{ex} %[2H5V27]
Từ một tờ giấy hình vuông, gấp đôi tờ giấy theo đường chéo ta được hình gì?
\choice
{\True Hình tam giác}
{Hình chữ nhật}
{Hình tứ giác}
{Hình tròn}
\loigiai{
Gấp đôi hình vuông theo đường chéo tạo thành hình tam giác vuông cân.
\textbf{Đáp án của câu hỏi là A}
}
\end{ex}

\begin{ex} %[2H5V27]
Cắt một hình vuông theo một đường thẳng nối hai điểm chính giữa của hai cạnh đối diện, ta được hai hình gì?
\choice
{Hai hình vuông}
{Hai hình tam giác}
{\True Hai hình chữ nhật bằng nhau}
{Hai hình tròn}
\loigiai{
Cắt đôi hình vuông theo trục đối xứng song song với cạnh ta được hai hình chữ nhật bằng nhau.
\textbf{Đáp án của câu hỏi là C}
}
\end{ex}

\begin{ex} %[2H5V27]
Cho hình chữ nhật, nếu cắt đi một góc của hình chữ nhật thì hình còn lại có thể có bao nhiêu cạnh?
\choice
{$3$ cạnh}
{\True $5$ cạnh (hình ngũ giác)}
{$2$ cạnh}
{$6$ cạnh}
\loigiai{
Cắt vát một góc của hình chữ nhật làm xuất hiện thêm 1 cạnh mới, tạo thành hình 5 cạnh (ngũ giác).
\textbf{Đáp án của câu hỏi là B}
}
\end{ex}

\begin{ex} %[2H5V27]
Ghép hai hình tam giác bằng nhau có thể tạo thành hình nào dưới đây?
\choice
{Hình tròn}
{Hình lục giác}
{Đường cong kín}
{\True Hình tứ giác}
\loigiai{
Ghép hai tam giác cùng chung một cạnh tạo thành một hình tứ giác.
\textbf{Đáp án của câu hỏi là D}
}
\end{ex}

\begin{ex} %[2H5V27]
Cho ba điểm $A, B, C$. Dùng thước kẻ đặt qua ba điểm thấy mép thước trùng khít với cả ba điểm. Ta kết luận
\choice
{\True Ba điểm $A, B, C$ thẳng hàng}
{Ba điểm tạo thành tam giác}
{Ba điểm không thẳng hàng}
{$B$ là điểm chính giữa}
\loigiai{
Mép thước thẳng đi qua cả ba điểm chứng tỏ ba điểm đó thẳng hàng.
\textbf{Đáp án của câu hỏi là A}
}
\end{ex}

\begin{ex} %[2H5V27]
Người ta xếp $4$ hình tam giác nhỏ bằng nhau thành một hình tam giác lớn. Số hình tam giác có trong hình lớn là
\choice
{$4$ hình tam giác}
{$6$ hình tam giác}
{\True $5$ hình tam giác}
{$8$ hình tam giác}
\loigiai{
Gồm 4 hình tam giác nhỏ và 1 hình tam giác lớn bao quanh: $4 + 1 = 5$ hình tam giác.
\textbf{Đáp án của câu hỏi là C}
}
\end{ex}

\begin{ex} %[2H5V27]
Dùng $6$ que tính dài bằng nhau xếp được nhiều nhất bao nhiêu hình tam giác rời nhau?
\choice
{$3$ hình tam giác}
{\True $2$ hình tam giác}
{$1$ hình tam giác}
{$4$ hình tam giác}
\loigiai{
Mỗi hình tam giác rời nhau cần 3 que tính: $6 : 3 = 2$ hình tam giác.
\textbf{Đáp án của câu hỏi là B}
}
\end{ex}

\begin{ex} %[2H5V27]
Để kiểm tra xem ba cái cọc có được cắm thẳng hàng hay không, người ta thường làm thế nào?
\choice
{Đo khoảng cách giữa các cọc}
{Dùng dây quấn quanh các cọc}
{Nhổ cọc lên kiểm tra}
{\True Đứng ở một đầu ngắm sao cho cọc thứ nhất che khuất cọc thứ hai và cọc thứ ba}
\loigiai{
Phương pháp ngắm thẳng hàng trong thực tế.
\textbf{Đáp án của câu hỏi là D}
}
\end{ex}

\begin{ex} %[2H5V27]
Một hình vuông được chia thành $4$ ô vuông nhỏ bằng nhau. Nếu tô màu $2$ ô vuông thì còn lại mấy ô vuông chưa tô màu?
\choice
{\True $2$ ô vuông}
{$1$ ô vuông}
{$3$ ô vuông}
{$4$ ô vuông}
\loigiai{
$4 - 2 = 2$ ô vuông.
\textbf{Đáp án của câu hỏi là A}
}
\end{ex}

\begin{ex} %[2H5V27]
Cắt một tờ giấy hình tam giác bằng một nhát kéo thẳng cắt ngang qua hai cạnh. Ta nhận được hai hình là
\choice
{Hai hình tam giác}
{Hai hình tứ giác}
{\True Một hình tam giác và một hình tứ giác}
{Một hình chữ nhật và một hình vuông}
\loigiai{
Phần trên là hình tam giác, phần dưới là hình tứ giác (hình thang).
\textbf{Đáp án của câu hỏi là C}
}
\end{ex}

\begin{ex} %[2H5V27]
Có bao nhiêu hình tứ giác trong một hình chữ nhật bị kẻ một đường thẳng chia đôi từ trên xuống dưới?
\choice
{$2$ hình tứ giác}
{\True $3$ hình tứ giác}
{$4$ hình tứ giác}
{$1$ hình tứ giác}
\loigiai{
Gồm 2 hình chữ nhật nhỏ và 1 hình chữ nhật lớn bao quanh: có tất cả 3 hình tứ giác.
\textbf{Đáp án của câu hỏi là B}
}
\end{ex}

\begin{ex} %[2H5V27]
Gấp một tờ giấy hình tròn làm tư rồi cắt theo đường cong viền. Khi mở ra ta nhận được
\choice
{Một đoạn thẳng}
{Một hình tam giác}
{Một đường gấp khúc}
{\True Hình có hoa văn đối xứng}
\loigiai{
Gấp và cắt giấy hình đối xứng.
\textbf{Đáp án của câu hỏi là D}
}
\end{ex}

\textbf{Phần 2. Trắc nghiệm đúng sai}
\setcounterpageref{ex}{0}

\begin{ex} %[2H5V27]
Xét tính đúng sai của các bài toán thực hành cắt ghép hình:
\choiceTF
{\True Hai hình tam giác bằng nhau luôn có thể ghép thành một hình tứ giác.}
{Một hình vuông cắt thành hai nửa thì hai nửa đó bắt buộc phải là hình tam giác.}
{\True Ba điểm A, B, C thẳng hàng nếu chúng cùng nằm trên một đường thẳng.}
{Dùng 4 que tính bằng nhau luôn ghép được một hình tam giác.}
\loigiai{
\begin{itemize}
    \item Ý a) Đúng: Ghép dọc theo một cạnh bằng nhau tạo thành hình tứ giác. Đúng.
    \item Ý b) Sai: Có thể cắt thành hai hình chữ nhật.
    \item Ý c) Đúng: Đúng theo định nghĩa.
    \item Ý d) Sai: Hình tam giác cần 3 que tính, 4 que tính ghép được hình vuông/tứ giác.
\end{itemize}
\begin{itemize}
    \item a) \textbf{Đúng}.
    \item b) \textbf{Sai}.
    \item c) \textbf{Đúng}.
    \item d) \textbf{Sai}.
\end{itemize}
}
\end{ex}

\begin{ex} %[2H5V27]
Người ta dùng các mảnh bìa hình vuông cạnh 1 cm để ghép thành hình chữ nhật dài 4 cm và rộng 2 cm:
\choiceTF
{\True Cần tất cả 8 mảnh bìa hình vuông nhỏ.}
{\True Chu vi hình chữ nhật ghép được là 12 cm.}
{Diện tích của hình chữ nhật là 12 cm vuông.}
{\True Nếu bớt đi 2 mảnh bìa ở một góc thì hình còn lại có 6 cạnh.}
\loigiai{
\begin{itemize}
    \item Ý a) Đúng: $4 \times 2 = 8$ mảnh bìa.
    \item Ý b) Đúng: $(4 + 2) \times 2 = 12\text{ cm}$.
    \item Ý c) Sai: Diện tích là $4 \times 2 = 8\text{ cm}^2$.
    \item Ý d) Đúng: Hình chữ L có 6 cạnh. Đúng.
\end{itemize}
\begin{itemize}
    \item a) \textbf{Đúng}.
    \item b) \textbf{Đúng}.
    \item c) \textbf{Sai}.
    \item d) \textbf{Đúng}.
\end{itemize}
}
\end{ex}

\textbf{Phần 3. Trả lời ngắn}
\setcounterpageref{ex}{0}

\begin{ex} %[2H5V27]
Có bao nhiêu hình tam giác có trong một hình tam giác lớn bị kẻ một đường thẳng từ đỉnh xuống đáy?
\shortans{3}
\loigiai{
\textbf{Đáp số:} $3$\\
Gồm 2 hình tam giác nhỏ và 1 hình tam giác lớn: $2 + 1 = 3$ hình.
}
\end{ex}

\begin{ex} %[2H5V27]
Cần ít nhất bao nhiêu que tính bằng nhau để xếp được 2 hình tam giác dính liền có chung 1 cạnh?
\shortans{5}
\loigiai{
\textbf{Đáp số:} $5$\\
Tam giác thứ nhất cần 3 que, tam giác thứ hai dùng chung 1 que nên cần thêm 2 que: $3 + 2 = 5$ que tính.
}
\end{ex}

\begin{ex} %[2H5V27]
Trong hình chữ nhật có kẻ hai đường chéo cắt nhau tại tâm. Hỏi có tất cả bao nhiêu hình tam giác?
\shortans{8}
\loigiai{
\textbf{Đáp số:} $8$\\
4 tam giác nhỏ đơn và 4 tam giác ghép đôi: $4 + 4 = 8$ hình tam giác.
}
\end{ex}

\begin{ex} %[2H5V27]
Cho 4 điểm trong đó không có 3 điểm nào thẳng hàng. Nối từng cặp điểm vẽ được bao nhiêu đoạn thẳng?
\shortans{6}
\loigiai{
\textbf{Đáp số:} $6$\\
Số đoạn thẳng: $4 \times 3 : 2 = 6$ đoạn thẳng.
}
\end{ex}

\begin{ex} %[2H5V27]
Một hình vuông có cạnh 10 cm. Chu vi của hình vuông đó bằng bao nhiêu xăng-ti-mét?
\shortans{40}
\loigiai{
\textbf{Đáp số:} $40$\\
$10 \times 4 = 40\text{ cm}$.
}
\end{ex}

\begin{ex} %[2H5V27]
Dùng 8 que tính bằng nhau để xếp thành một hình vuông. Hỏi mỗi cạnh của hình vuông có mấy que tính?
\shortans{2}
\loigiai{
\textbf{Đáp số:} $2$\\
$8 : 4 = 2$ que tính.
}
\end{ex}
